\documentclass[11pt]{article}
\usepackage[utf8]{inputenc}
\usepackage[T1]{fontenc}
\usepackage[margin=1in]{geometry}
\usepackage{amsmath, amssymb}
\usepackage{booktabs}
\usepackage{array}
\usepackage{hyperref}

\begin{document}

% =============================
% Exam definition (included)
% =============================
\begin{center}
{\LARGE First Exam — Lessons 0 \& 1 (Version A)}\\[0.75em]
\end{center}

\noindent\textbf{Time allowed:} 60 minutes\\
\textbf{Resources:} Printed course notes and personal handwritten notes only. Electronic devices, calculators, and statistical tables are not permitted.\\
\textbf{Instructions:} Answer all four questions. Show algebraic steps and clearly label each part. Unless stated otherwise, provide exact values (fractions, radicals, exponentials). Partial credit is available when reasoning is clearly explained.

\vspace{0.75em}
\noindent\begin{tabular}{@{}p{2.2cm}p{8.2cm}r@{}}
\toprule
\textbf{Question} & \textbf{Topic} & \textbf{Points}\\
\midrule
1 & Probability structures and independence & 25\\
2 & Bayes' rule and conditional reasoning & 25\\
3 & Exponential model derivations & 25\\
4 & Interpreting summaries and convergence & 25\\
\midrule
\textbf{Total} & & \textbf{100}\\
\bottomrule
\end{tabular}

\vspace{1em}
\hrule
\vspace{1em}

\section*{Question 1 — Probability Foundations (25 points)}

A coin is tossed twice. Throughout parts (a)–(c), assume the coin is fair.

\begin{enumerate}
  \item (6 pts) Specify a probability space $\,(\Omega, \mathcal{F}, P)\,$ for this experiment. Clearly list the sample space, a natural $\sigma$-algebra, and the probability measure.
  \item (6 pts) Let $A$ be the event ``at least one head appears'' and $B$ the event ``exactly one head appears.'' Compute $P(A)$, $P(B)$, and $P(A \cap B)$.
  \item (5 pts) Are $A$ and $B$ independent? Justify your answer using the formal definition.
  \item (8 pts) Now suppose the coin has probability $p$ of landing heads, where $0 < p < 1$. With the same definitions of $A$ and $B$, derive expressions for $P_p(A)$, $P_p(B)$, and $P_p(A \cap B)$. For which value(s) of $p$ are $A$ and $B$ independent? Provide algebraic reasoning.
\end{enumerate}

\vspace{0.75em}
\hrule
\vspace{0.75em}

\section*{Question 2 — Bayes' Rule in Component Testing (25 points)}

A factory purchases circuit boards from two suppliers. Boards from $S_A$ account for $\tfrac{3}{5}$ of the total, and boards from $S_B$ account for the remaining $\tfrac{2}{5}$. Historical quality data show that $\tfrac{1}{10}$ of boards from $S_A$ are defective, while $\tfrac{1}{20}$ of boards from $S_B$ are defective. Each board is screened by a test that returns ``positive'' for a defective board with probability $\tfrac{9}{10}$, and returns ``positive'' for a non-defective board with probability $\tfrac{1}{20}$.

\medskip
\noindent\textit{Recall:} sensitivity $= P(\text{Positive} \mid \text{Defective})$ and specificity $= P(\text{Negative} \mid \text{Non-defective})$. The false positive rate is $1-\text{specificity} = P(\text{Positive} \mid \text{Non-defective})$.

\begin{enumerate}
  \item (7 pts) Compute the probability that a randomly selected board produces a positive test result.
  \item (8 pts) Given a positive test result, determine the probability that the board is defective. Express your answer as a reduced fraction.
  \item (6 pts) Given a negative test result, determine the probability that the board is defective. Express your answer as a reduced fraction.
  \item (4 pts) Let $\pi = P(\text{Defective})$, $s = P(\text{Positive} \mid \text{Defective})$, and $f = P(\text{Positive} \mid \text{Non-defective})$. Derive a formula for $P(\text{Defective} \mid \text{Positive})$ in terms of $\pi$, $s$, and $f$, and briefly state how decreasing $f$ affects this probability.
\end{enumerate}

\vspace{0.75em}
\hrule
\vspace{0.75em}

\section*{Question 3 — Exponential Waiting-Time Model (25 points)}

Let $T$ be the waiting time (in hours) until the next service request arrives, and suppose $T$ follows an exponential distribution with rate $\lambda = \tfrac{1}{4}$.

\medskip
\noindent\textbf{Recall:} The probability density function of an exponential random variable with rate $\lambda$ is given by
\[
f_T(t) =
\begin{cases}
\lambda e^{-\lambda t}, & t \ge 0,\\
0, & t < 0.
\end{cases}
\]

\begin{enumerate}
  \item (5 pts) Derive the cumulative distribution function $F_T(t)$.
  \item (5 pts) Using your result from part (1), compute $P(1 \le T \le 4)$ and leave the answer in exponential form.
  \item (9 pts) Compute $E[T]$ and $\operatorname{Var}(T)$ directly from integrals. Show the essential integration steps leading to exact values.
  \item (6 pts) Define $Y = 3T + 2$. Compute $E[Y]$ and $\operatorname{Var}(Y)$ using properties of expectation and variance. Provide exact values.
\end{enumerate}

\vspace{0.75em}
\hrule
\vspace{0.75em}

\section*{Question 4 — Interpreting Summaries and Convergence (25 points)}

Table~\ref{tab:a-sim-sol} summarises the results of a simulation: for each sample size $n$, 1{,}000 independent samples of size $n$ were drawn from an exponential distribution with mean $4$, and the sample mean was recorded for each run.

\begin{table}[h]
\centering
\caption{Behaviour of Sample Means}\label{tab:a-sim-sol}
\begin{tabular}{@{}rcc@{}}
\toprule
\textbf{Sample size $n$} & \textbf{Average of the 1,000 sample means} & \textbf{SD of the 1,000 sample means}\\
\midrule
5  & 4.08 & 1.80 \\
20 & 4.02 & 0.90 \\
80 & 4.01 & 0.44 \\
\bottomrule
\end{tabular}
\end{table}

\begin{table}[h]
\centering
\caption{Selected QQ-Plot Quantiles}\label{tab:a-qq-sol}
\begin{tabular}{@{}rc@{}}
\toprule
\textbf{Normal quantile} & \textbf{Sample quantile}\\
\midrule
-2.0 & -3.1 \\
-1.0 & -1.6 \\
 0.0 &  0.2 \\
 1.0 &  1.9 \\
 2.0 &  3.4 \\
\bottomrule
\end{tabular}
\end{table}

\newpage

% =============================
% Solutions and grading notes
% =============================
\begin{center}
{\LARGE Solutions and Grading Notes (Version A)}\\[0.5em]
\end{center}

\section*{Question 1 — Probability Foundations (25 points)}

\begin{enumerate}
  \item \textbf{Probability space (6 pts)}\\
  Sample space: $\Omega = \{\mathrm{HH},\, \mathrm{HT},\, \mathrm{TH},\, \mathrm{TT}\}$.\\
  Natural $\sigma$-algebra: $\mathcal{F} = 2^{\Omega}$.\\
  Probability measure: For the fair coin, $P(\omega) = \tfrac{1}{4}$ for each $\omega \in \Omega$.\\
  Partial credit: correct $\Omega$ (2 pts), acceptable $\mathcal{F}$ (2 pts), correct probabilities summing to 1 (2 pts).

  \item \textbf{Event probabilities (6 pts)}\\
  $A = \{\mathrm{HH},\, \mathrm{HT},\, \mathrm{TH}\}$ so $P(A) = \tfrac{3}{4}$.\\
  $B = \{\mathrm{HT},\, \mathrm{TH}\}$ so $P(B) = \tfrac{1}{2}$.\\
  $A \cap B = \{\mathrm{HT},\, \mathrm{TH}\}$ so $P(A \cap B) = \tfrac{1}{2}$.\\
  Award 3 pts for correct event descriptions, 3 pts for the resulting probabilities.

  \item \textbf{Independence check (5 pts)}\\
  $P(A)P(B) = \tfrac{3}{4} \cdot \tfrac{1}{2} = \tfrac{3}{8} \neq \tfrac{1}{2} = P(A \cap B)$. Therefore $A$ and $B$ are \emph{not} independent.\\
  Credit: 3 pts for the product comparison, 2 pts for the explicit conclusion.

  \item \textbf{Biased coin analysis (8 pts)}\\
  With head probability $p$, $P_p(\mathrm{HH}) = p^2$, $P_p(\mathrm{HT}) = p(1-p)$, $P_p(\mathrm{TH}) = (1-p)p$, $P_p(\mathrm{TT}) = (1-p)^2$.\\
  $P_p(A) = 1 - (1-p)^2 = 2p - p^2$, \quad $P_p(B) = 2p(1-p)$, \quad $P_p(A \cap B) = 2p(1-p)$.\\
  Independence would require $2p(1-p) = (2p - p^2)\,2p(1-p)$. If $0<p<1$, division by $2p(1-p)$ yields $1 = 2p - p^2$, i.e. $(p-1)^2 = 0$. The only solution is $p=1$, which lies outside $(0,1)$ and renders $B$ impossible. Hence for $0<p<1$, $A$ and $B$ are not independent.\\
  Award 5 pts for correct expressions, 3 pts for the independence argument.
\end{enumerate}

\section*{Question 2 — Bayes' Rule in Component Testing (25 points)}

Let $D$ denote ``defective'' and $+$ denote ``positive test.''

\begin{enumerate}
  \item \textbf{Overall positive probability (7 pts)}\\
  Prevalence: $P(D) = \tfrac{3}{5}\cdot\tfrac{1}{10} + \tfrac{2}{5}\cdot\tfrac{1}{20} = \tfrac{3}{50} + \tfrac{1}{50} = \tfrac{2}{25}$.\\
  Therefore $P(+) = P(+\mid D)P(D) + P(+\mid D^c)P(D^c) = \tfrac{9}{10}\cdot\tfrac{2}{25} + \tfrac{1}{20}\cdot\tfrac{23}{25} = \tfrac{36}{500} + \tfrac{23}{500} = \tfrac{59}{500}$.

  \item \textbf{Positive predictive value (8 pts)}\\
  $P(D\mid +) = \dfrac{P(+\mid D)P(D)}{P(+)} = \dfrac{\tfrac{9}{10}\cdot\tfrac{2}{25}}{\tfrac{59}{500}} = \dfrac{\tfrac{9}{125}}{\tfrac{59}{500}} = \dfrac{36}{59}$.

  \item \textbf{Posterior after a negative (6 pts)}\\
  $P(-\mid D) = 1 - \tfrac{9}{10} = \tfrac{1}{10}$, \; $P(-\mid D^c) = 1 - \tfrac{1}{20} = \tfrac{19}{20}$.\\
  $P(-) = \tfrac{1}{10}\cdot\tfrac{2}{25} + \tfrac{19}{20}\cdot\tfrac{23}{25} = \tfrac{1}{125} + \tfrac{437}{500} = \tfrac{441}{500}$.\\
  Hence $P(D\mid -) = \dfrac{\tfrac{1}{10}\cdot\tfrac{2}{25}}{\tfrac{441}{500}} = \dfrac{\tfrac{1}{125}}{\tfrac{441}{500}} = \dfrac{4}{441}$.

  \item \textbf{General formula and interpretation (4 pts)}\\
  Bayes' rule gives $P(D\mid +) = \dfrac{s\pi}{s\pi + f(1-\pi)}$. Decreasing $f$ (the false-positive rate) reduces the denominator while leaving the numerator unchanged, thereby increasing $P(D\mid +)$.
\end{enumerate}

\section*{Question 3 — Exponential Waiting-Time Model (25 points)}

\begin{enumerate}
  \item \textbf{CDF (5 pts)}\\
  For $\lambda = \tfrac{1}{4}$, $\displaystyle F_T(t) = \begin{cases} 0, & t<0,\\ 1 - e^{-t/4}, & t\ge 0.\end{cases}$

  \item \textbf{Probability of an interval (5 pts)}\\
  $P(1 \le T \le 4) = F_T(4) - F_T(1) = (1 - e^{-1}) - (1 - e^{-1/4}) = e^{-1/4} - e^{-1}$.

  \item \textbf{Expectation and variance via integration (9 pts)}\\
  $E[T] = \int_0^\infty t \cdot \tfrac{1}{4} e^{-t/4}\, dt = 4$.\\
  $E[T^2] = \int_0^\infty t^2 \cdot \tfrac{1}{4} e^{-t/4}\, dt = 16\int_0^\infty x^2 e^{-x}dx = 32$.\\
  $\operatorname{Var}(T) = 32 - 16 = 16$.

  \item \textbf{Affine transformation (6 pts)}\\
  $E[Y] = E[3T+2] = 3E[T] + 2 = 14$, \quad $\operatorname{Var}(Y) = 3^2\operatorname{Var}(T) = 144$.
\end{enumerate}

\section*{Question 4 — Interpreting Summaries and Convergence (25 points)}

\begin{enumerate}
  \item \textbf{Law of Large Numbers interpretation (9 pts)}\\
  The average of the sample means in Table~\ref{tab:a-sim-sol} stays close to the true mean 4: 4.08 for $n=5$, 4.02 for $n=20$, and 4.01 for $n=80$. As $n$ increases, the averages move closer to 4 and fluctuate less, illustrating that sample means converge toward the population mean, as predicted by the Law of Large Numbers.

  \item \textbf{Central Limit Theorem interpretation (8 pts)}\\
  The standard deviation of the 1{,}000 sample means decreases roughly by a factor close to $1/\sqrt{n}$: 1.80 for $n=5$, 0.90 for $n=20$, and 0.44 for $n=80$. This shrinking variability reflects the CLT prediction that $\bar X_n$ becomes more concentrated and approximately normal with variance $\sigma^2/n$ as $n$ grows. Full credit requires explicitly linking the observed reductions to the CLT scaling.

  \item \textbf{QQ-plot assessment and recommendation (8 pts)}\\
  Sample quantiles are more extreme in both tails than the normal quantiles (e.g., $-3.1<-2$ and $3.4>2$), indicating a QQ-plot that bends away from the 45° line with heavier tails. Therefore the dataset is unlikely to follow a standard normal distribution closely. Suggested actions: adopt a heavier-tailed model (e.g., Student-$t$), employ robust procedures, or transform the data before applying normal-based methods. Full credit needs description of the pattern, a clear normality conclusion, and an appropriate recommendation.
\end{enumerate}

\end{document}
